\paragraph{更新设计与效果归因。}
实验比较的对象是完整的Block3块级更新：共享advantage、联合重要性比率与按有效长度加权的block归一化。
第\ref{sec:method}节分别刻画了这一更新的跨token信用、梯度尺度与裁剪几何。
因此，实测性能差异对应完整更新，而非单独的advantage平均。
组件贡献是另一层实验问题，可通过固定ratio与归一化、仅改变反馈共享来研究。
信号估计分析在明确假设下描述平均的性质，任务准确率则由benchmark评测衡量。

\paragraph{统计单位与权重报告。}
每个报告的实验组使用一个训练seed。配对题目区间描述给定训练权重下的评测不确定性，
而非不同训练seed之间的波动；区间为逐点估计，未作多重比较校正。
AIME每组30题，多解决一题对应Pass@8增加3.33个百分点。
除Step200外，我们同时保留各保存点的表现轨迹；Step200是回顾性采用的统一计划内终点，
并非预注册的权重选择规则。窗口与加权变体来自探索性实验，独立运行可进一步评估其差异的可重复性。

\paragraph{协议与任务覆盖。}
历史结果按各自的prompt、seed处理、教师及实现设置组织，每一组对应其协议内的方法比较。
两组完整留存对照共用DAPO衍生训练池；重叠审计检查留存数据，不覆盖语义等价或上游模型数据暴露。
4B completion比较两种训练方法，未包含匹配的原始学生评测；官方指令版额外包含该对照。
独立数据集、训练seed与师生对构成进一步拓展评测的自然维度。

\paragraph{互补诊断。}
提示probe检查优化前的行为，留存的训练轨迹描述后续演化。
熵图刻画当前策略访问的前缀，boxed缺失率衡量一种特定格式属性。
梯度范数与长度停止率记录数值及生成行为，符号分歧衡量反馈在位置间的重新分配，
benchmark准确率则提供独立的任务评价。
联合解释这些指标有助于区分输入适配、优化动态与数学正确性，而不以其中一项替代其余评价。

\paragraph{基线与计算范围。}
Sampled-token OPD是匹配实验基线，其他OPD方案在相关工作中讨论。
Block3改变已有教师反馈驱动更新的方式，仍对响应的每个有效位置评分。
因此，当前评测关注这一更新的行为与准确率，而非减少教师评分工作量。
逐任务、逐保存点报告反映了实测表现的变化；后续可在共同预算下比较组件选择及准确率与成本的权衡。
